% Autor: Elias Welt
\section{Benutzeroberfläche und Design (UI/UX)}
\subsection{Styling-Strategie}
Das Projekt verwendet natives CSS mit CSS-Variablen (Custom Properties) für das Theming. Dies bietet maximale Flexibilität ohne den Overhead großer Frameworks.
Globale Styles sind in \texttt{index.css} definiert, während komponentenspezifische Styles (z. B. \texttt{Chat.css}) lokal importiert werden.

\subsection{Dark Mode Implementierung}
Der Dunkelmodus ist nicht nur eine kosmetische Spielerei, sondern eine zentrale Anforderung. Er wurde über CSS-Variablen realisiert, die je nach Attribut am HTML-Tag ihre Werte ändern.

\begin{lstlisting}[language=CSS, caption={Theme-Switching mittels CSS-Variablen}]
:root {
  /* Standard (Dark Mode) */
  --primary-color: #222831;
  --text-color: #FFFFFF;
}

[data-theme='light'] {
  /* Light Mode Overrides */
  --primary-color: #FFFFFF;
  --text-color: #222831;
}

body {
  background-color: var(--primary-color);
  color: var(--text-color);
}
\end{lstlisting}

Ein \texttt{useEffect}-Hook in der \texttt{App.jsx} prüft beim Start den \texttt{localStorage}, ob der Nutzer eine Präferenz gespeichert hat, und setzt das Attribut entsprechend. Dies verhindert den "Flash of Incorrect Theme" (FOUC).
